% Part: applied-modal-logics
% Chapter: epistemic-logic
% Section: language-epistemic-logic

\documentclass[../../../include/open-logic-section]{subfiles}

\begin{document}

\olfileid[hi]{aml}{el}{lan}

\olsection{ज्ञानात्मक तर्कशास्त्र की भाषा}

\begin{defn}
मान लें कि $G$ कर्ता-संकेतों का समुच्चय है। बहु-कर्ता ज्ञानात्मक तर्कशास्त्र
की मूल भाषा में निम्नलिखित होते हैं:
\begin{enumerate}
  \tagitem{prvFalse}{!!{falsity} का प्रतिज्ञप्तिक अचर~$\lfalse$।}{}
  \tagitem{prvTrue}{!!{truth} का प्रतिज्ञप्तिक अचर~$\ltrue$।}{}
  \item !!{propositional variable}ों का !!{denumerable} समुच्चय: $\Obj
    p_0$, $\Obj p_1$, $\Obj p_2$, \dots
  \item प्रतिज्ञप्तिक संयोजक: \startycommalist
  \iftag{prvNot}{\ycomma $\lnot$ (निषेध)}{}%
  \iftag{prvAnd}{\ycomma $\land$ (संयोजन)}{}%
  \iftag{prvOr}{\ycomma $\lor$ (वियोजन)}{}%
  \iftag{prvIf}{\ycomma $\lif$ (!!{conditional})}{}%
  \iftag{prvIff}{\ycomma $\liff$ (!!{biconditional})}{}।
  \item प्रत्येक $a \in G$ के लिए ज्ञान-संकारक $\Knows_a$।
\end{enumerate}
\end{defn}

यदि हमारी रुचि अपने तंत्र में केवल एक कर्ता के ज्ञान में है, तो हम समुच्चय
~$G$ और अलग-अलग कर्ताओं का उल्लेख छोड़ सकते हैं। उस स्थिति में हमारे पास
केवल मूल संकारक~$\Knows$ होता है।

\begin{defn}
ज्ञानात्मक भाषा की \emph{!!{formula}ें} आगमनात्मक रूप से इस प्रकार परिभाषित
होती हैं:
\begin{enumerate}
\tagitem{prvFalse}{$\lfalse$ एक परमाण्विक !!{formula} है।}{}

\tagitem{prvTrue}{$\ltrue$ एक परमाण्विक !!{formula} है।}{}

\item प्रत्येक प्रतिज्ञप्तिक चर $\Obj p_i$ एक (परमाण्विक) !!{formula} है।

\tagitem{prvNot}{यदि $!A$ एक !!{formula} है, तो $\lnot !A$ एक
  !!{formula} है।}{}

\tagitem{prvAnd}{यदि $!A$ और $!B$ !!{formula} हैं, तो $(!A \land
  !B)$ एक !!{formula} है।}{}

\tagitem{prvOr}{यदि $!A$ और $!B$ !!{formula} हैं, तो $(!A \lor !B)$
  एक !!{formula} है।}{}

\tagitem{prvIf}{यदि $!A$ और $!B$ !!{formula} हैं, तो $(!A \lif !B)$
  एक !!{formula} है।}{}

\tagitem{prvIff}{यदि $!A$ और $!B$ !!{formula} हैं, तो $(!A \liff !B)$
  एक !!{formula} है।}{}

\item यदि $!A$ एक !!{formula} और $a \in G$ है, तो $\Knows_a !A$ एक
  !!{formula} है।

\tagitem{limitClause}{इसके अतिरिक्त कुछ भी !!{formula} नहीं है।}{}
\end{enumerate}
\end{defn}

यदि किसी !!{formula}~$!A$ में $\Knows_a$ नहीं आता, तो हम उसे
\emph{मोडल-मुक्त} कहते हैं।

\begin{defn}
  संकारक $\Knows$ का उद्देश्य व्यक्तिगत ज्ञान को संकेतित करना है, जबकि
  $\EKnows$, जिसे प्रायः ``हर कोई जानता है'' पढ़ते हैं, सामूहिक ज्ञान को
  संकेतित करता है। जहाँ $G' \subseteq G$, वहाँ हम $\EKnows_{G'} !A$ को
  निम्न का संक्षिप्त रूप मानते हैं:
  \[\bigwedge_{b \in G'} \Knows_b !A.\]
\end{defn}

हम ज्ञान की एक और भी प्रबल धारणा, अर्थात कर्ताओं के समूह~$G$ में
\emph{सार्वज्ञान}, भी परिभाषित कर सकते हैं। जब कोई सूचना कर्ताओं के समूह
में सार्वज्ञान होती है, तो इसका अर्थ है कि उस समूह के कर्ताओं के प्रत्येक
संयोजन के लिए वे सभी जानते हैं कि दूसरे जानते हैं कि दूसरे जानते हैं
\dots और यह क्रम अनंत तक चलता है। यह सामूहिक ज्ञान से बहुत अधिक प्रबल है,
और ऐसे संबंधात्मक प्रतिरूप आसानी से बनाए जा सकते हैं जिनमें कोई !!{formula}
सामूहिक ज्ञान तो हो, किंतु सार्वज्ञान न हो। ``समूह~$G$ में यह सार्वज्ञान है
कि~$!A$'' को हम $\CKnows_G !A$ से संकेतित करेंगे।

\end{document}
